\chapter{Ecuaciones diofánticas}
\label{md:diofanticas}

Una ecuación diofántica busca soluciones \emph{enteras}.
Encontrar una solución real no resuelve el problema, y cualquier
restricción de positividad debe imponerse además de la integridad.

\section{Ecuación lineal en dos incógnitas}
\label{md:diofantica-lineal}

Sean $a,b,c\in\Z$, $(a,b)\ne(0,0)$ y $d=\mcd(a,b)$.
La ecuación
\[
 ax+by=c
\]
tiene soluciones enteras si y solo si $d\mid c$.
La necesidad se sigue de que $d$ divide cualquier combinación;
la suficiencia usa Bézout: si $as+bt=d$, una solución es
$x_0=sc/d$, $y_0=tc/d$.
Todas las soluciones son
\[
 \boxed{x=x_0+\frac bd\,k,\qquad y=y_0-\frac ad\,k,
 \quad k\in\Z.}
\]
Al restar dos soluciones, la coprimalidad de $a/d,b/d$ obliga a que
sus diferencias tengan esos múltiplos. La fórmula también cubre
el caso en que uno de los coeficientes sea cero.
Si $a=b=0$, hay todas las parejas enteras cuando $c=0$, y ninguna si $c\ne0$.

Para resolver: calcular $d$, comprobar $d\mid c$, obtener una solución
por Euclides extendido y escribir la familia completa. Si se exigen
$x,y\ge0$ o $x,y>0$, se sustituyen las expresiones y se determina
el intervalo de valores enteros de $k$, respetando el signo de cada coeficiente.

\begin{ejemplo}
En $14x+21y=7$, una solución es $(-1,1)$ y $d=7$.
Todas son $x=-1+3k$, $y=1-2k$.
No hay soluciones con ambas coordenadas no negativas: la primera
exigiría $k\ge1$ y la segunda $k\le0$.
\end{ejemplo}

\section{Diferencia de cuadrados y factorización}

La ecuación $x^2-y^2=n$ se transforma en
\[
 (x-y)(x+y)=n.
\]
Una pareja de factores $uv=n$ proporciona enteros
$x=(u+v)/2$, $y=(v-u)/2$ exactamente cuando $u,v$ tienen la misma
paridad. Para $x>y>0$ se toman $0<u<v$.

La \textbf{factorización de Fermat}, para $n>1$ impar, busca
$n=x^2-y^2$. Se empieza con $x=\lceil\sqrt n\rceil$ y se incrementa
$x$ hasta que $x^2-n$ sea un cuadrado; entonces
$n=(x-y)(x+y)$.
Existe siempre la descomposición trivial
\[
 n=\left(\frac{n+1}{2}\right)^2-
   \left(\frac{n-1}{2}\right)^2.
\]
Si no aparece una factorización antes de ese valor de $x$, $n$ es primo.
El método resulta especialmente adecuado cuando los factores están próximos.

\section{Ternas pitagóricas y descenso}

Una \textbf{terna pitagórica} es una solución positiva de
$x^2+y^2=z^2$. Es \textbf{primitiva} si $\mcd(x,y,z)=1$;
en ese caso los tres números son coprimos dos a dos, un cateto es par
y el otro y la hipotenusa son impares.

Todas las ternas primitivas, intercambiando los catetos si es necesario, son
\[
 x=s^2-t^2,\qquad y=2st,\qquad z=s^2+t^2,
\]
donde $s>t>0$, $\mcd(s,t)=1$ y $s,t$ tienen distinta paridad.
La idea es factorizar $y^2=(z-x)(z+x)$: después de dividir por $4$,
se obtiene un producto de coprimos que debe ser cuadrado en cada factor.
Las ternas positivas generales se obtienen multiplicando una primitiva
por un entero positivo.

El \textbf{descenso infinito} parte de una supuesta solución positiva
y construye otra con una cantidad positiva estrictamente menor.
Buena ordenación impide repetir ese descenso: basta elegir primero
una solución mínima. Así se demuestra que $x^4+y^4=z^2$ no tiene
soluciones positivas.

El \textbf{último teorema de Fermat} afirma que
$x^n+y^n=z^n$ no tiene soluciones enteras positivas para $n>2$.
El caso de exponente $4$ se deduce del resultado anterior;
no se extiende ese argumento por sí solo a todos los exponentes.
